\documentclass[10pt,a4paper]{amsart}
\usepackage[margin=19mm]{geometry}
\usepackage{amsmath,amssymb,mathtools}
\usepackage[scr=rsfs,cal=euler]{mathalfa}
\usepackage{xfrac}
\usepackage[unicode,hidelinks,bookmarks=false]{hyperref}
\newcommand{\ie}{i.e.\ }
\newcommand{\cf}{cf.\ }
\newcommand{\resp}{respectively\ }
\newcommand{\Subsection}{\subsection}
\providecommand{\nobreakdash}{\nobreak\hbox{-}}
\setlength{\emergencystretch}{2em}
\multlinegap0pt
\usepackage{bm}
\usepackage{bbm}
\usepackage{dsfont}
\usepackage{xspace}
\usepackage{cancel}
\usepackage{mathtools}
\usepackage{amsthm}
\makeatletter
\@ifundefined{thm}{\newtheorem{thm}{Theorem}}{}
\makeatother
\newcommand{\I}{{\bf I}}
\newcommand{\cD}{\underline{\bf D}}
\newcommand{\cL}{\underline{\bf L}}
\newcommand{\cU}{\underline{\bf U}}
\newcommand{\cV}{\underline{\bf V}}
\newcommand{\cW}{\underline{\bf W}}
\newcommand{\cX}{\underline{\bf X}}
\newcommand{\cY}{\underline{\bf Y}}
\newcommand{\cZ}{\underline{\bf Z}}
\title[A note on generalized tensor CUR approximation for tensor pa]{A note on generalized tensor CUR approximation for tensor pairs and tensor triplets based on the tubal product}
\author{Salman Ahmadi-Asl, Naeim Rezaeian, Keivan Ramazani}
\date{Source: arXiv:2305.00754v4 (CC BY 4.0)}
\begin{document}
\maketitle

\noindent\emph{Source and licence.} A note on generalized tensor CUR approximation for tensor pairs and tensor triplets based on the tubal product. Version arXiv:2305.00754v4 (CC BY 4.0), distributed under the Creative Commons Attribution 4.0 International licence, as recorded in the arXiv licence metadata for this exact version and shown on the abstract page https://arxiv.org/abs/2305.00754v4. Full rights notice: licenses/arxiv-2305.00754-CC-BY-4.0.txt.\par\smallskip

\noindent\textit{Editorial scope.} Counted unit: the Thm whose source label is \texttt{None} in the article (source0.tex lines 553--562, source label \texttt{None}), delivered together with the article's own complete original proof of it (source0.tex lines 564--583, one \texttt{proof} environment of 20 lines). No mathematics is written, repaired or re-derived: statement and proof are the article's own text line for line. Required context retained uncounted: rsvd (lines 546--548). Every definition, display and label of those lines is retained unchanged. Omitted from this package and not redistributed: 1--545, 549--552, 563--563, 584--691. Nothing inside a retained excerpt is omitted or altered: every retained body is delivered exactly as the article's lines are written, so no omission receipt of any kind accompanies this package. Citations: the retained excerpts carry no citation command at all, so no bibliography entry of the article is transcribed and no reference is redistributed. Reference adaptation: none was needed and none was made; every reference inside the delivered unit resolves inside it, so no unresolved reference or citation is produced. Printed slips kept literal: no printed word, symbol, hypothesis or number of the counted statement or of its proof was altered. Layout and class: the article's own class is \texttt{elsarticle} with packages including graphicx, amssymb, algcompatible, algpseudocode, multirow, xspace, hyperref, listings, url, csquotes, euscript, amsfonts, epsfig, enumerate; the wrapper is the lane's pinned \texttt{amsart} editorial wrapper at 10pt on A4, which restates the theorem-like environments the retained text needs and adds the article's own macro definitions that the retained text needs (\texttt{\textbackslash I}, \texttt{\textbackslash cD}, \texttt{\textbackslash cL}, \texttt{\textbackslash cU}, \texttt{\textbackslash cV}, \texttt{\textbackslash cW}, \texttt{\textbackslash cX}, \texttt{\textbackslash cY}, \texttt{\textbackslash cZ}), with extra wrapper packages (bm, bbm, dsfont, xspace, cancel, mathtools). The delivered numbering is the unit's own. Assets: no retained span contains a figure environment, a graphics-inclusion command, a PDF-page inclusion, a table or a TikZ picture; no image file is copied into this package, the package declares no asset dependency and no image file is redistributed with it. Licence: version arXiv:2305.00754v4 of this article is distributed under the Creative Commons Attribution 4.0 International licence, as recorded in the arXiv licence metadata for this exact version and shown on the abstract page https://arxiv.org/abs/2305.00754v4. A full rights notice is delivered at licenses/arxiv-2305.00754-CC-BY-4.0.txt. Characters: every retained excerpt is pure ASCII, so the delivered engine, XeLaTeX with its default Latin Modern text font, reports no missing character.
\begin{equation}\label{rsvd}
\cX=\cL*\cD_1*\cW^T,\quad \cY=\cL*\cD_2*\cU^T,\quad \cZ=\cV*\cD_3*\cW^T,
\end{equation}

\begin{thm} Let $\cX\in\mathbb{R}^{I_1\times I_2\times I_3}$, $\cY\in\mathbb{R}^{I_1\times I_4\times I_3}$, and $\cZ\in\mathbb{R}^{I_5\times I_2\times I_3}$ be given, then

\begin{itemize}
\item a) If $\cY$ and $\cZ$ are nonsingular tensors, then the selected lateral and horizontal slice indices from the TCUR approximation of $\cY^{-1}*\cX*\cZ^{-1}$ are identical to the index vectors ${\bf p}_1$ and ${\bf s}_1$, respectively,  obtained from a GTCUR approximation of the tensor triples $(\cX,\cY,\cZ)$.

\item  b) In the particular case where $\cY = \underline{\I}$ and $\cZ = \underline{\I}$, the GTCUR approximation of $(\cX,\underline{\bf I},\underline{\bf I})$ coincides with the TCUR approximation of $\cX$.

\item c) For a special choice of $\cY=\underline{\bf I}$, a GTCUR approximation of $(\cX,\underline{\bf I},\cZ)$ coincides with the GTCUR approximation of $(\cX,\cZ)$. Similar results can be stated for $\cZ = \underline{\bf I}$.
\end{itemize}
\end{thm}

\begin{proof}
a) From \ref{rsvd}, it is not difficult to see that $\cY^{-1}=\cU*{\cD}_2^{-1}*\cL^{-1},\,\cZ^{-1}=\cW*\cD_3^{-1}*\cV^T$, so we have 
\begin{eqnarray}
\nonumber
\cY^{-1}*\cX*\cZ^{-1}=(\cU*{\cD}_2^{-1}*\cL^{-1})*(\cL*\cD_1*\cW^T)*(\cW*\cD_3^{-1}*\cV^T)=\\\label{prf}
\cU*(\cD_2^{-1}\cD_1\cD_3^{-1})*\cV^T.
\end{eqnarray}
This means that the selected lateral and horizontal slice indices of $\cY^{-1}*\cX*\cZ^{-1}$ using the TCUR approximation are the same as the index vectors ${\bf p}_1$ and ${\bf s}_1$, respectively, obtained from a GTCUR approximation of $(\cX,\cY,\cZ).$\\

b) If $\cY=\underline{\I}$ and $\cZ=\underline{\I}$, then $\cL=\cU*{\cD}_2^{-1}$ and $\cW^T=\cD_3^{-1}*\cV^T$. Substituting them in the first part of \eqref{rsvd}, we have $\cX=\cU*(\cD_2^{-1}*\cD_1*\cD_3^{-1})*\cV^T$. This indicates that both algorithms utilize the same tensors $\cU$ and $\cV$  for selecting the horizontal and lateral slices, respectively. This completes the proof.


c) If $\cY=\underline{\I}$, then from the second part of \eqref{rsvd}, we have $\cL=\cU*\cD_2^{-1}$. If we substitute this in the first part of \eqref{rsvd}, we arrive at 
\begin{eqnarray}\label{rs_1}
\cX&=&\cU*(\cD_2^{-1}*\cD_1)*\cW^T,\\\label{rs_2}
\cZ&=&\cV*\cD_3*\cW^T.  
\end{eqnarray}
It is evident that equations \eqref{rs_1} and \eqref{rs_2} demonstrate that the GTCUR approximation of the tensor pairs $(\cX,\cZ)$ yields identical indices for both horizontal and lateral slice sampling as the GTCUR approximation of the tensor triples $(\cX,\underline{\bf I},\cZ)$. This completes the proof.

\end{proof}

\end{document}
